\documentclass[12pt, letterpaper]{article}

% --- PREÁMBULO --- 
\usepackage[utf8]{inputenc} 
\usepackage[T1]{fontenc} 
\usepackage[spanish]{babel} 
\usepackage{xcolor} % Para colorear las respuestas

% Márgenes de 2cm en todos los lados 
\usepackage[letterpaper, top=2cm, bottom=2cm, left=2cm, right=2cm]{geometry} 
\usepackage{graphicx} 
\usepackage[most]{tcolorbox} 
\usepackage{enumitem} 
\usepackage{amssymb} 
\usepackage{amsmath}
\usepackage{multicol}

% --- PAQUETE PARA FONDO ---
\usepackage{eso-pic}


% --- FUENTE --- 
\usepackage{helvet} 
\renewcommand{\familydefault}{\sfdefault}

% --- COMANDO PARA CASILLAS DE RESPUESTA PEQUEÑAS --- 
\newcommand{\boxfill}[1][1.5cm]{\tcbox[on line, size=minimal, colframe=black, colback=white, boxrule=0.5pt, arc=1mm, boxsep=3pt, left=2pt, right=2pt, top=2pt, bottom=2pt]{\hspace*{#1}\vphantom{M}}}

% =================================================================
% CONFIGURACIÓN PARA MOSTRAR / OCULTAR RESPUESTAS
% =================================================================
\newif\ifshowanswers

% ---> MODO PROFESOR (Muestra respuestas en azul): Deja la siguiente línea sin comentar.
% ---> MODO ESTUDIANTE (Deja las casillas en blanco para imprimir): Pon un % al inicio de la línea.
%\showanswerstrue 
\showanswersfalse 

% Comandos auxiliares para pintar las respuestas
\newcommand{\casillaVerdadera}{\ifshowanswers \textcolor{blue}{$\blacksquare$} \else $\square$ \fi}
\newcommand{\casillaFalsa}{\ifshowanswers $\square$ \else $\square$ \fi}
\newcommand{\casillaNum}[1]{\ifshowanswers \textcolor{blue}{\textbf{#1}} \else $\square$ \fi}

\begin{document}
	
	% --- ENCABEZADO --- 
	\begin{center} 
		\textbf{TALLER DE REPASO PRE-EXAMEN}\\[0.1cm] 
		TAU 6: ¿Se puede ver la energía? 
	\end{center}
	
	% --- PUNTO 1 --- 
	\noindent \textbf{1.} En física, ¿cuál de las siguientes opciones define mejor la palabra ``energía''? Marca con una \textbf{X} la opción correcta. 
	
	\vspace{0.3cm} 
	\begin{itemize}[itemsep=0.4cm] 
		\item[ \casillaFalsa ] Es una sustancia mágica de color azul que flota en el aire. 
		\item[ \casillaFalsa ] Es una fuerza que se gasta y desaparece completamente cuando encendemos una máquina.
		\item[ \casillaVerdadera ] Es la capacidad invisible que tiene algo para efectuar o realizar un trabajo. 
	\end{itemize}
	
	\vspace{0.5cm}
	
	% --- PUNTO 2 --- 
	\noindent \textbf{2.} Lee cada afirmación sobre las características de la energía y marca con una \textbf{X} si es Verdadera (V) o Falsa (F): 
	
	\vspace{0.3cm} 
	\begin{itemize}[itemsep=0.6cm, label={}] 
		\item \textbf{A.} Podemos ver la energía flotando en el aire si usamos unas gafas \\ científicas o un telescopio muy potente. \hfill  \casillaFalsa  V \quad  \casillaVerdadera  F 
		\item \textbf{B.} Según la Ley de Causalidad, si vemos un carro moverse (el efecto),\\ sabemos que hay una energía invisible trabajando (la causa). \hfill  \casillaVerdadera  V \quad  \casillaFalsa  F 
	\end{itemize}
	
	\vspace{0.5cm}
	
	% --- PUNTO 3 --- 
	\noindent \textbf{3.} La energía se manifiesta de muchas formas. Relaciona el "apellido" de la energía con estos \textbf{nuevos ejemplos} poniendo la letra en el paréntesis correspondiente: 
	
	\vspace{0.3cm}
	\begin{multicols}{2}
		\begin{itemize}[itemsep=0.5cm, label={}]
			\item \textbf{A.} Energía Lúminica
			\item \textbf{B.} Energía Química
			\item \textbf{C.} Energía Calórica
			\item \textbf{D.} Energía Eléctrica
		\end{itemize}
		\columnbreak
		\begin{itemize}[itemsep=0.5cm, label={}]
			\item ( \ifshowanswers \textcolor{blue}{\textbf{D}} \else \hspace{0.5cm} \fi ) En los rayos de una gran tormenta.
			\item ( \ifshowanswers \textcolor{blue}{\textbf{A}} \else \hspace{0.5cm} \fi ) En la pantalla encendida de un celular.
			\item ( \ifshowanswers \textcolor{blue}{\textbf{C}} \else \hspace{0.5cm} \fi ) En una fogata para calentarnos.
			\item ( \ifshowanswers \textcolor{blue}{\textbf{B}} \else \hspace{0.5cm} \fi ) El carbón o el desayuno que comes.
		\end{itemize}
	\end{multicols} 
	
	\newpage
	% --- PUNTO 4 --- 
	\vspace*{0.5cm} 
	\noindent \textbf{4.} En la mitología, Proteo era un personaje muy escurridizo. ¿Cuál era su poder especial y por qué el libro de ciencias lo compara con la energía? 
	
	\vspace{0.3cm} 
	\begin{itemize}[itemsep=0.5cm] 
		\item[ \casillaFalsa ] Su poder era tener una fuerza infinita para mover montañas, como la electricidad. 
		\item[ \casillaFalsa ] Su poder era hacerse invisible para siempre, igual que la energía cuando se destruye. 
		\item[ \casillaVerdadera ] Su poder era cambiar de apariencia y transformarse (en agua, dragón, etc.), tal como la energía cambia de forma constantemente. 
	\end{itemize}
	
	\vspace{0.6cm}
	
	% --- PUNTO 5 --- 
	\noindent \textbf{5.} Si un compañero te dice que la ``energía lumínica'' y la ``energía eléctrica'' son sustancias totalmente distintas y separadas en el universo, ¿qué le responderías según lo que leíste en el TAU? 
	
	\vspace{0.4cm} 
	\begin{itemize}[itemsep=0.5cm] 
		\item[ \casillaVerdadera ] "¡Te equivocas! Es una única energía; esas palabras son solo 'apellidos' para saber cómo se está transformando."
		\item[ \casillaFalsa ] "¡Tienes razón! Cada vez que encendemos una máquina se crea una energía nueva de la nada." 
		\item[ \casillaFalsa ] "¡Depende! Solo la energía eléctrica es verdadera, las demás son solo inventos de la ciencia." 
	\end{itemize}
	
	\vspace{1.5cm}
	

	
\end{document}